\section{Introduction}
\label{sec:introduction}

Commercial large language models (LLMs) can produce high-quality curricula and
support instructional systems design (ISD) from needs analysis through assessment
~\cite{isdagentbench}. Institutional use nevertheless brings recurring inference
cost, external-API dependence, and third-party transfer of learner or institutional
data. On-premise open-weight small language models (SLMs) offer more control, but as
the model shrinks, the harness that manages state, coordinates calls, and checks
outputs must carry more of the system-level burden.

Existing ISD and lesson-planning agents emphasize automation, personalization, or
quality with capable foundation models. They do not treat preservation of
instructional-design quality during model downscaling as a joint design and
evaluation objective. A deployment-oriented study must establish both quality at
each feasible size and whether the relative gap changes as capacity falls.

A second gap concerns constructive alignment among learning objectives, activities,
and assessments~\cite{biggs1996enhancing,biggs2011teaching}. ISD benchmarks assess
this property mainly through LLM-judge rubrics~\cite{isdagentbench}, placing the
measurement instrument in the same broad family as the generator. Porter/Webb
indices and embedding methods can quantify correspondence to human-authored
standards~\cite{porter2002measuring,webb1997criteria}, but no ISD benchmark uses
such methods to examine the intrinsic alignment of a wholly machine-generated
design, where both sides of every relation are system outputs.

Our response begins from the deployment constraint. Each model call receives one
bounded ISD concern rather than the full open-ended, multi-constraint task. The
premise that narrower subproblems can be more tractable is inherited from work on
decomposed prompting and small-model orchestration
~\cite{zhou2023leasttomost,khot2023decomposed,ferraz2024decrim,
shen2024multillmagent,belcak2025slm}; it is not tested here because every ablation
retains decomposition. Nor do we assume that its benefit grows as models shrink:
harness effects can be non-monotonic across model tiers~\cite{cho2026harness}, so
the scale interaction is an empirical question.

Decomposition creates its own problem: separately generated artifacts must retain
one pedagogical intent. AlignmentGraph-ISD therefore uses constructive alignment to
shape a shared graph of objectives, activities, assessments, and their relations.
Specialists read only relevant graph context and update scoped state; a verifier
checks the evolving design, and a graph dump exposes the final relations for audit.
The graph is a substrate for constrained generation, not a claim that prose is
aligned by construction (\Cref{sec:method}).

The evaluation separates quality level from scale interaction. RQ1a compares the
harness with six agents at each Qwen size; RQ1b tests whether those gaps widen as
the model shrinks. The \texttt{addie\_median}, \texttt{trajectory\_score}, and
\texttt{total\_score} judge signals receive identical analyses and equal reporting
status. ADDIE leads the narrative only because its artifact rubric matches RQ1's
content-quality construct. RQ2 asks how measurable signals of constructive
alignment vary across the same ladder, without claiming that automated scores
validate the construct itself (\Cref{sec:rqs,sec:protocol}).

This paper makes two contributions:
\begin{enumerate}
  \item \textbf{AlignmentGraph-ISD}, a multi-agent ADDIE harness whose alignment
  graph is explicit, checkable state and an auditable educator-facing artifact.
  \item \textbf{An automated constructive-alignment evaluation protocol} extending
  ISD-Agent-Bench with a no-LLM, non-hierarchical panel of seven signals from
  mechanically independent textual-correspondence and cognitive-demand instrument
  families. It operationalizes established criteria rather than proposing a new
  construct~\cite{biggs2011teaching,qm2018rubric,porter2002measuring,
  webb1997criteria}; cross-family agreement is an initial validity bridge, while
  expert-human validation remains future work.
\end{enumerate}

The experiment spans four Qwen sizes from 0.8B to 9B, three independent runs per
cell, and three ablations inside the ladder, with between-run variation reported
for every estimate. Following this Introduction, \Cref{sec:related-work} locates
the two gaps; \Cref{sec:method} presents the harness and evaluation protocol;
\Cref{sec:experimental-methodology,sec:results} describe the experimental design
and findings; \Cref{sec:discussion} interprets the results and bounds their
validity; and \Cref{sec:conclusion} closes the paper.
